\section{Randomized --- NPPoly}
%%% generated by the blueprint pipeline from TCSlib.Complexity.Randomized.NPPoly
\subsection*{Randomized --- NPPoly: Introduction}

Declarations of the Lean module \texttt{TCSlib.Complexity.Randomized.NPPoly}.

\subsection*{Randomized --- NPPoly: Results}

\begin{theorem}[A Good Shared Randomness String]
\lean{Randomized.exists_good_coins}
\label{Randomized.exists_good_coins}
\leanok
\uses{Randomized.blockCount, Randomized.blockCount_add_not, Randomized.exists_forall_not_of_randProb_le, Randomized.one_sub_pow_le_half_pow, Randomized.polyLen, Randomized.randProb, Randomized.randProb_bool_false, Randomized.randProb_congr, Randomized.randProb_tail_le, Turing.pairEncode}
\difficulty{6}
Let $L$ and $C$ be languages of bit strings, let $f$ map bit strings to bit strings, and
let $a,k$ be nonnegative integers. Suppose that for every bit string $x$, a uniformly
random string $r$ of length $a(|x|+1)^k$ satisfies $f(\langle x,r\rangle)\in C$ if and
only if $x\in L$ with probability at least $2/3$, where $\langle x,r\rangle$ is the
pairing encoding described above. For every nonnegative integer $n$, put $M=36(n+1)$ and
$q=a(n+1)^k$. There is a bit string $\rho$ of length $Mq$ such that, for every bit
string $v$ of length $n$, membership $v\in L$ is equivalent to more than half of the $M$
strings formed by taking the $q$-bit blocks of $\rho$ (starting with block $i$ for each
$0\le i<M$) making $f(\langle v,\text{block}_i\rangle)$ belong to $C$.
\end{theorem}

\begin{definition}[Polynomial-size nondeterministic circuit recognition]
\lean{Language.InNPPoly}
\statementsource{Yap83}{pdf-1a201b302069-p003-b002}
\statementsource{Vad12}{
  pdf-bf2ce5422088-p057-b006,
  pdf-bf2ce5422088-p058-b002
}
\statementsource{arora-barak-ch07-randomized-computation}{pdf-26d95080dbde-p018-b002}
\label{Language.InNPPoly}
\leanok
\uses{BoolCircuit.DAGCircuit, BoolCircuit.DAGCircuit.IsFaninTwo, BoolCircuit.DAGCircuit.eval, BoolCircuit.DAGCircuit.size}
A language $L$ of Boolean strings belongs to the class described here if there are
natural numbers $a,k$, a function $m:\mathbb{N}\to\mathbb{N}$, and a Boolean circuit
$C_n$ for each input length $n$, with $n+m(n)$ input bits, such that each circuit has
fan-in at most two and size at most $a(n+1)^k$. For every Boolean string $x$, membership
$x\in L$ is equivalent to the existence of a witness of length $m(|x|)$ for which
$C_{|x|}$ outputs true on the concatenation of $x$ and that witness.
\end{definition}

\begin{definition}[Boolean values on wires]
\lean{Randomized.wireVal}
\label{Randomized.wireVal}
\leanok
For a natural number $N$ and a Boolean valuation $v$ on the wires indexed by
$0,\ldots,N-1$, extend the valuation to the disjoint union of those wires with one
additional Boolean wire by assigning $v$ on the first summand and the wire's own Boolean
value on the second.
\end{definition}

\begin{theorem}[List of evaluated wire values]
\lean{Randomized.ofFn_inputValues_vectorOfList}
\label{Randomized.ofFn_inputValues_vectorOfList}
\leanok
\uses{BoolCircuit.inputValues, BoolCircuit.vectorOfList, BoolCircuit.vectorOfList_ofFn, Randomized.wireVal}
\difficulty{4}
Let $w$ be a finite list of wires, each either an input position in $	ext{Fin}(N)$ or a
Boolean constant, and let $v$ assign a Boolean value to each input position. The list
obtained by evaluating each wire in $w$—using $v$ for input positions and the given
value for constants—is the image of $w$ under this evaluation.
\end{theorem}

\begin{theorem}[Polynomial verification from polynomial-size wirings]
\lean{Randomized.inNPPoly_of_wiring}
\label{Randomized.inNPPoly_of_wiring}
\leanok
\uses{BoolCircuit.DAGCircuit.bufferInputs, BoolCircuit.DAGCircuit.bufferInputs_eval, BoolCircuit.DAGCircuit.bufferInputs_structure, BoolCircuit.DAGCircuit.size, BoolCircuit.DAGCircuitFamily.mem_language_ofFn, BoolCircuit.vectorOfList, Complexity.P, Complexity.P_subset_PPoly, Language.InNPPoly, Language.InPPoly, Randomized.ofFn_inputValues_vectorOfList, Randomized.wireVal}
\difficulty{5}
Let $L$ and $W$ be languages of bit strings, with $W$ decidable in polynomial time. Let
$m:\mathbb{N}\to\mathbb{N}$, and for each $n$ let a wiring be a list whose entries are
either a bit constant or one of the $n+m(n)$ input bits. Suppose there are natural
numbers $b,j$ such that, for every $n$, the wiring length plus one and the total number
$n+m(n)$ of input bits are each at most $b(n+1)^j$. If $x\in L$ exactly when some
assignment to $m(|x|)$ witness bits makes the wiring output a string in $W$, then $L$
has polynomial-size fan-in-two circuit verifiers with polynomially many witness bits.
\end{theorem}

\begin{theorem}[Length after duplicating list entries]
\lean{Randomized.length_flatMap_dup}
\label{Randomized.length_flatMap_dup}
\leanok
\difficulty{4}
For any type $\alpha$ and any finite list $l$ of elements of $\alpha$, replacing each
entry by two consecutive copies gives a list of length $2|l|$.
\end{theorem}

\begin{theorem}[Duplicating elements under flat map]
\lean{Randomized.flatMap_dup_map}
\label{Randomized.flatMap_dup_map}
\leanok
\difficulty{4}
For any function $g$ from a type $\alpha$ to a type $\beta$ and any finite list $l$ of
elements of $\alpha$, mapping each element to two consecutive copies of its image and
concatenating gives the same list as first mapping $g$ over $l$ and then duplicating
each resulting element.
\end{theorem}

\begin{theorem}[Length of concatenated list blocks]
\lean{Randomized.length_flatMap_range}
\label{Randomized.length_flatMap_range}
\leanok
\difficulty{4}
Let $g$ assign a list to each natural number, and let $K,c$ be natural numbers. If
$g(j)$ has length $c$ for every $j<K$, then concatenating the lists $g(0),\ldots,g(K-1)$
produces a list of length $Kc$.
\end{theorem}

\begin{theorem}[Extracting a block from a concatenation]
\lean{Randomized.flatMap_range_drop_take}
\label{Randomized.flatMap_range_drop_take}
\leanok
\uses{Randomized.length_flatMap_range}
\difficulty{4}
Let $g$ assign a finite list to each natural number, and fix natural numbers $c$, $K$,
and $i$. If $g(j)$ has length $c$ for every $j<K$ and $i<K$, then the length-$c$ block
beginning at position $ic$ in the concatenation $g(0),g(1),\ldots,g(K-1)$ is exactly
$g(i)$.
\end{theorem}

\begin{theorem}[Length of a list block]
\lean{Randomized.length_block}
\label{Randomized.length_block}
\leanok
\difficulty{3}
Let $l$ be a list and let $K,c,i$ be natural numbers. If $l$ has length $Kc$ and $i<K$,
then the block obtained by dropping the first $ic$ elements of $l$ and taking the next
$c$ elements has length $c$.
\end{theorem}

\begin{theorem}[Randomized reductions to 3-SAT have polynomial witnesses]
\lean{Randomized.BPNP_subset_NPPoly}
\label{Randomized.BPNP_subset_NPPoly}
\leanok
\uses{Complexity.P, Complexity.PolyTimeComputable, Complexity.PolyTimeComputable.append, Complexity.PolyTimeComputable.comp, Complexity.PolyTimeComputable.output_length_le, Complexity.SAT3, Complexity.SAT3_mem_NP, Complexity.blockAt, Complexity.mem_P_of_blockMajority, Complexity.pairFstD, Complexity.pairFstD_pairEncode, Complexity.pairSndD, Complexity.pairSndD_pairEncode, Complexity.polyTimeComputable_pairFstD, Complexity.polyTimeComputable_pairSndD, Complexity.preimage_mem_P, Language.InNPPoly, Randomized.BPNP, Randomized.composeMap, Randomized.composeMap_pairEncode, Randomized.exists_good_coins, Randomized.exists_ofFn_eq, Randomized.flatMap_dup_map, Randomized.flatMap_range_drop_take, Randomized.inNPPoly_of_wiring, Randomized.length_block, Randomized.length_flatMap_dup, Randomized.length_flatMap_range, Randomized.polyLen, Randomized.polyLen_length_pairEncode_le, Randomized.polyTimeComputable_composeMap, Randomized.wireVal, Turing.MultiTapeTM.indicator, Turing.pairEncode}
\difficulty{7}
For every verifier model $E$, every language that randomized polynomial-time reduces to
the language of satisfiable CNF formulas of width at most $3$ has polynomially bounded
witnesses verifiable by a family of fan-in-two Boolean circuits of polynomial size:
membership is equivalent to the existence of a witness that makes the appropriate
circuit accept.
\end{theorem}

\begin{theorem}[NP Circuits with Polynomial Size]
\lean{Randomized.NP_subset_NPPoly}
\label{Randomized.NP_subset_NPPoly}
\leanok
\uses{Complexity.NP, Language.InNPPoly, Randomized.BPNP_subset_NPPoly, Randomized.NP_subset_BPNP}
\difficulty{1}
Every language in $\mathrm{NP}$ has a family of fan-in-two Boolean circuits with
polynomially bounded size and polynomially many witness input bits, such that an input
belongs to the language exactly when some assignment to its witness bits makes the
corresponding circuit output true.
\end{theorem}

\begin{theorem}[A coNP inclusion from randomized SAT]
\lean{Randomized.coNP_subset_NPPoly_of_compl_SAT3_mem_BPNP}
\label{Randomized.coNP_subset_NPPoly_of_compl_SAT3_mem_BPNP}
\leanok
\uses{Complexity.SAT3, Complexity.SAT3_NPHard, Complexity.coNP, Language.InNPPoly, Randomized.BPNP, Randomized.BPNP_subset_NPPoly, Randomized.RandReducible.of_polyTimeReducible_left}
\difficulty{4}
If the complement of the language of satisfiable 3-CNF formulas belongs to BPNP, then
every language in coNP has polynomial-size, fan-in-two circuit certificates; that is,
coNP is contained in the class of languages recognized by such circuits with
polynomially many existential Boolean inputs.
\end{theorem}
